\begin{abstract}
\noindent Can a learner that acquires inference rules come to reason soundly, and with what justification? We consider a learner that imitates human inferences, is then conditioned on coherence (no good arguments for both $P$ and $\neg P$), and receives occasional feedback from the world. We study this question for formal mathematics, informal mathematics, and the checking of physics solutions. All guarantees are relative to stated assumptions, chiefly that the target is realizable in a structured hypothesis class and that data and designations are truthful.

The answer is a division of labour among the learning signals, each with provable strengths and a provable residue.
\begin{itemize}
\item \textbf{Search forces worst-case soundness.} Accuracy on the distribution of human steps says nothing about a reasoner that searches: one tonk-like rule of arbitrarily small probability mass makes everything derivable. Under realizability, cautious acceptance over schematic hypotheses is sound against every adaptive prover. Its price is a combinatorial dimension of the hypothesis class.
\item \textbf{Imitation identifies rules.} For rules cited by name, positive examples identify a schematic calculus exactly at coupon-collector rates, and this generalizes to derivations of any size. But systematic human errors are indistinguishable from rules.
\item \textbf{Coherence is eliminative negative data.} It pins the target down exactly when the target has no coherent uniform proper extension. Classical logic has none, by Post-completeness, and is pinned by a single coherence datum. Coherence is blind among intermediate logics and Popperian for arithmetic. Carnap's categoricity problem can be read, in part, as Gold's problem in the dual space of valuations.
\item \textbf{Formal mathematics.} We prove an end-to-end theorem for a two-tier learner: audit, then assert, with bold conjectures in a sandbox. It holds under realizability, a frequency floor, truthful designation and an exact refutation oracle, and each of several ingredients is shown to be necessary.
\item \textbf{Informal mathematics.} We model it as the shadow of a latent formalization. A bounded-gap verifier is sound against adaptive provers when the formalization is realizable. With complete data, validity is identified up to informal-validity equivalence and no finer. Counterexample objects can be exponentially more informative than paradoxes.
\item \textbf{Physics.} Contexts receive a filter semantics. No truth-functional eternalist reading tracks in-context truth, and exports from idealized contexts need certificates of finite information radius. A structured-solution checker is sound against adversarial solvers relative to a judgment layer, consisting of the reading of the problem and the adequacy of the top model, which we prove ineliminable. We illustrate it on one olympiad problem.
\end{itemize}
We also answer, in part, two questions raised in the notes that motivated this work: whether a steeper simplicity penalty leads from imitation to norms, and whether there is a ``philosophical completeness theorem''.

Every main result was attacked by independent referees and repaired where needed. Core discrete results are machine-checked in Lean~4, and small symbolic experiments illustrate the phenomena. We argue that principled justification beyond proof has a definite shape: proof within contexts, certified bridges between them, and an explicit, minimized, ineliminable judgment layer, together with learning signals whose one-sided limits are known.
\end{abstract}
